\chapter[Band 35: Rechtskürzbare Halbgruppen]{Band 35 auf einen Blick}
\noindent\textbf{Rechtskürzbare Halbgruppen}\par\medskip
Vorausgesetzt sind Halbgruppen und die natürliche Addition.
Links- und Rechtskürzbarkeit werden als getrennte Eigenschaften geführt;
in Band 36 werden beide gemeinsam verlangt.

\begin{BandUebersicht}
\textbf{Kürzungsaxiomatik}
Zur Halbgruppenstruktur kommt die angegebene Kürzungsrichtung:\UebFormel{\begin{aligned}
&a,b,c\in A,\ b\star a=c\star a\ \vdash b=c.
\end{aligned}}Das Strukturprädikat lautet \(\RightCancellativeSemigroupAlg A\star\).
\UebQuellen{\FormulaRefAuto{RightCancellativeSemigroupDef}[def]}
&
\textbf{Die natürliche Addition}
\UebFormel{\begin{aligned}
&\RightCancellativeSemigroupAlg{\mathbb N}{+},\\
&a,b,c\in\mathbb N,\ b+a=c+a\ \vdash b=c.
\end{aligned}}
\UebQuellen{\FormulaRefAuto{NaturalAdditionRightCancellativeSemigroup}[thm];
\FormulaRefAuto{NaturalAdditionRightCancellative}[thm]} \\
\addlinespace[.8ex]

\textbf{Kommutativität und Linkskürzung}
\UebFormel{\begin{aligned}
&\CommutativeSemigroupAlg A\star,\qquad\\
&\LeftCancellativeSemigroupAlg A\star.
\end{aligned}}
\UebQuellen{\FormulaRefAuto{CommutativeSemigroupDef}[def];
\FormulaRefAuto{LeftCancellativeSemigroupDef}[def]}
&
\textbf{Daraus folgt Rechtskürzung}
\UebFormel{\begin{aligned}
&\CommutativeSemigroupAlg A\star,\\
&\LeftCancellativeSemigroupAlg A\star\ \vdash\\
&\qquad\RightCancellativeSemigroupAlg A\star.
\end{aligned}}
\UebQuellen{\FormulaRefAuto{CommutativeLeftCancellativeSemigroupImpliesRightCancellativeSemigroup}[thm]} \\
\addlinespace[.8ex]

\textbf{Homomorphismen und Isomorphismen}
Für zwei Strukturen derselben Klasse bedeutet ein Homomorphismus
\(H(f)\):\UebFormel{\begin{aligned}
&f:A\to B,\\
&x,y\in A\ \vdash\\
&\quad f(x\star y)=f(x)\diamond f(y).
\end{aligned}}Ein Isomorphismus \(I(f)\) ist ein bijektiver Homomorphismus:\UebFormel{\begin{aligned}
&I(f)\ \leftrightarrow H(f)\land f:A\bij B.
\end{aligned}}
\UebQuellen{\FormulaRefAuto{RightCancellativeSemigroupHomDef}[def];
\FormulaRefAuto{RightCancellativeSemigroupIsoDef}[def]}
&
\textbf{Komposition erhält die Struktur}
Für Homomorphismen \(f:A\to B\), \(g:B\to C\) zwischen
Strukturen dieser Klasse ist \(g\circ f\) wieder ein Homomorphismus.
Mit \(\bullet\) als Operation auf \(C\) gilt für \(x,y\in A\):\UebFormel{\begin{aligned}
&(g\circ f)(x\star y)\\
&\quad=(g\circ f)(x)\mathbin{\bullet}(g\circ f)(y).
\end{aligned}}Sind \(f,g\) Isomorphismen, so ist auch \(g\circ f\) ein Isomorphismus.
\UebQuellen{\FormulaRefAuto{SemigroupHomComposition}[thm];
\FormulaRefAuto{SemigroupIsoComposition}[thm]} \\
\end{BandUebersicht}

